\section{How Infinity Could Emerge}

If the universe is, or tends toward, infinite, how could that come to be, and what would it even be like? This section draws out the imaginative heart of the eternal story from the two cosmic stories, carefully. We keep the question open. The point is to make infinity thinkable, not to decide it.

Two stances frame everything that follows, and they are worth separating before any mechanism is named. \emph{Potential} infinity is a process that never finishes, finite at every stage, unbounded only in the limit it heads toward. \emph{Actual} infinity is a completed totality, all of it present at once, with no further stage to add. The first is a verb. The second is a noun. Almost every move in this section pushes infinity toward the first reading and away from the second.

\subsection{Seven ways infinity could emerge}

There is not one way for infinity to be. There are several, and they apply to different kinds of infinity, extent, duration, divisibility, and cardinality. A theory can take one route for time and a different route for space. The seven below are not rivals so much as a toolbox.

\begin{enumerate}
\item \textbf{As an unbounded process, potential infinity.} Infinity is never created, it is approached. The knit runs forever, finite at every beat, and infinity is just the name for the open horizon it tends toward. What is real is only the finite stages.
\item \textbf{As given and eternal, actual infinity.} Infinity is not created at all, it simply is, complete, with no origin. The cost is the classic paradoxes of a finished infinite.
\item \textbf{As a boundary, compactification.} Infinity is added to a finite-looking space as one tidy edge, the limit all outward paths share. You do not travel to it. This is literally the boundary at infinity of the bulk.
\item \textbf{As the continuum from the discrete, coarse-graining.} A discrete set has no points between its elements. Make it dense and complete it, and you get a continuum, infinitely many points in any interval. The mesh stays discrete at the base. The continuum is the emergent large-scale description.
\item \textbf{As exponential growth, vastness fast.} Negative curvature opens space up so quickly it becomes effectively unbounded almost at once. This is the bridge from a tiny seed to an effectively infinite extent, exactly the bulk's exponential roominess.
\item \textbf{From nothing, via creation and the arrow.} A finite beginning, then unbounded growth driven by the emergent arrow. The infinity is the future the process heads toward, not a thing that was made.
\item \textbf{As the totality of possibilities, ensembles.} The infinite is not in one world but in the collection of all of them.
\end{enumerate}

\begin{table}[ht]
\centering
\begin{tabular}{@{}lll@{}}
\toprule
route & infinity is & kind \\
\midrule
unbounded process & a never-completed limit & potential \\
given and eternal & a complete totality & actual \\
boundary & one shared outer edge & boundary object \\
continuum & density between discrete points & emergent \\
exponential growth & vastness reached fast & effective \\
creation and arrow & the open future & potential \\
ensembles & all worlds together & actual (of worlds) \\
\bottomrule
\end{tabular}
\caption{Seven routes to infinity, and what each one makes infinity be.}
\label{tab:infinity-routes}
\end{table}

\begin{result}[The synthesis]
Infinity is pushed to the edges. It is the limit of a process, a boundary object, or an emergent description. Rarely is it a completed actual thing sitting in the middle of reality.
\end{result}

\subsection{Finite-but-growing, the cleanest how}

The crispest version is the fourth and sixth routes combined, sharpened to a single picture.

At any single beat the universe has finitely many docks. But over all beats the count grows without bound. The wake adds new docks at the open frontier, beat after beat, and nothing caps the total.

So the universe is never actually infinite. It only tends toward infinity, and that tending never finishes. Infinity is the open-ended horizon of the growth, not a place the universe arrives at.

\begin{definition}[Finite-but-growing]
A flow is \emph{finite-but-growing} when the mesh holds finitely many docks at every beat, while the number of docks increases without bound across beats. Infinity is the limit of the count, never a value the count attains.
\end{definition}

Whether the completed limit is real or only a way of speaking is the metaphysical fork. The mathematics is identical either way. The difference is one of interpretation, not of any quantity you could measure inside the flow.

\subsection{Infinity as the default, it might need no cause}

There is a stronger possibility. That infinity is the default. An eternal, beginningless, edgeless mesh, the natural state, with no first beat to explain.

The reason this is even coherent is reversibility. The knit is collide then stream, and both phases are reversible. Started from a generic state the mesh churns forever, activity steady, no decay, and the whole evolution is exactly reversible. So every state has a unique past, and the history extends backward forever. No first beat is needed.

\begin{theorem}[Reversibility removes the first beat]
If the knit is exactly reversible, every state of the flow has a unique predecessor under it. The backward chain of predecessors never terminates. An eternal, beginningless history is therefore self-consistent, and no first beat is required.
\end{theorem}

On this reading infinity may even be the more natural option, for four reasons.

\begin{itemize}
\item \textbf{No privileged point.} An infinite homogeneous mesh has no special origin. A finite one has a center or seed that demands explanation. Infinity removes the why-here-why-then question.
\item \textbf{No boundary.} A finite universe needs an edge or an initial condition. An infinite eternal one needs neither. It just is, maximally symmetric.
\item \textbf{A fixed point.} An infinite eternal churn is already-everything, the natural attractor, where the growing state is always incomplete.
\item \textbf{Potential, not completed.} Infinite available expanse can mean unbounded potential rather than a completed actual infinity. That slides back toward finite-but-growing with no first beat.
\end{itemize}

Crucially this stays computable. Actual infinite extent is fine for a local knit, as long as every finite region is finitely computable. You never compute the whole at once. What is still avoided is a completed infinity as a single graspable object, and the real continuum. So an infinite past, or an infinite mesh, is allowed. An actual completed totality is not.

\begin{principle}[The computability line]
Infinite extent and infinite past are permitted, because the knit is local and every finite region resolves in finite work. What is forbidden is a completed infinite totality grasped at once, and the uncountable continuum as a base object. Locality keeps the infinite computable by never asking for all of it together.
\end{principle}

\subsection{The spectrum, finite to infinite is a ladder}

Finite or infinite is too coarse a cut. It is a spectrum, and a theory can be finite on one axis and infinite on another, across extent, duration, divisibility, cardinality, and information. The positions below are rungs on that ladder.

\begin{table}[ht]
\centering
\begin{tabular}{@{}ll@{}}
\toprule
position & claim \\
\midrule
finitism & only finite objects exist, no actual infinity of any kind \\
finite and bounded & finite, closed, no edge, like a sphere's surface \\
\textbf{finite but growing} & finite at every beat, unbounded over time (this substrate) \\
potentially infinite & reality is an unbounded process, infinity its never-completed limit \\
actually infinite (extent) & space is infinite in size now \\
actually infinite (continuum) & the real continuum is a completed totality, uncountable \\
infinite ensembles & infinitely many worlds, or all structures, exist \\
\bottomrule
\end{tabular}
\caption{The ladder from strict finitism up to infinite ensembles. This substrate sits at finite-but-growing.}
\label{tab:infinity-ladder}
\end{table}

The drift is downward-and-back. Away from infinity is actual and fundamental, and toward infinity is potential, a boundary, or emergent. The base is finite and discrete. The continuum coarse-grains out of it, and lives only at large scale.

\subsection{What it would feel like}

This is the part to sit with. Each item below is a felt distinction, the lived texture of a claim that is easy to state and hard to picture.

\begin{itemize}
\item \textbf{Counting is the model.} A process can be unbounded and always finite at once. There is no largest number you will ever reach, yet every number you reach is finite. No-largest does not mean an infinite one exists, it means for each, a bigger one. That is the felt sense of an endlessly growing universe, always finite, never finished.
\item \textbf{No edge does not mean infinite.} Walk any direction on a sphere's surface and you come back. No wall, no boundary, no outside, yet finite. A universe can be edgeless and still finite. Edgeless and endless are different feelings.
\item \textbf{A verb, not a noun.} The universe is growing without end (potential) is not the same claim as the universe contains infinitely many stars right now (actual). The first feels like a verb, a becoming. The second like a noun, a completed thing. The difference is completion, not size.
\item \textbf{Infinity as a horizon, not a place.} The boundary at infinity is a horizon of directions, a limit you approach and a screen, not a faraway spot with room in it.
\item \textbf{The eternal churn, from inside.} An eternal, beginningless, edgeless mesh feels maximally symmetric and nothing-special. The knit has simply always been running, every beat with a unique past behind it and a unique future ahead. A fire that was never lit and never winds down.
\item \textbf{Reality accretes.} On the growing reading, new beats are added at the leading edge and stay. The now is the most recent slice of an ever-growing whole. Time is lived duration, the future genuinely open, novelty real. The many become one and are increased by one.
\end{itemize}

\subsection{Where this substrate lands, and why it stays open}

The working default is \textbf{finite-but-growing}. A hyperbolic crystal finite at every beat, unbounded over all beats, discrete at the base with the continuum emergent. The finite-at-every-beat property belongs to the wake, the growing edge that adds docks one frontier at a time.

On that reading infinity is allowed only two non-fundamental roles. It is the boundary the growth approaches but never reaches, and it is the emergent continuum you see when you stop resolving docks. Everything real at every beat is finite.

But the eternal-beginningless-infinite picture is fully self-consistent too. The reversible knit needs no first beat. So it remains a live open axiom, exactly the fork left to the reader in the two cosmic stories.

\begin{principle}[Reversibility is the discriminator]
A reversible knit needs no beginning. So the choice between finite-but-growing and eternal-and-infinite is not forced by the mechanics. Both doors stay open, and reversibility is what keeps the second one unlocked.
\end{principle}

The mechanisms are settled. Each of the seven routes is a clean piece of mathematics or a property of the knit, and finite-but-growing is the substrate's default reading. Which one the universe actually realizes is not settled, and this section does not pretend to settle it. The how is in hand. The which is left open on purpose.
